% --- START OF FILE 07_compose.tex ---

\chapter{Jetpack Compose: UI Moderna e Dichiarativa}

\section{Da XML a Jetpack Compose}
Storicamente, Android usava un approccio \textbf{Imperativo} per la UI: si creava un albero di View in XML e si modificava il loro stato tramite codice Kotlin/Java (usando \texttt{findViewById}, \texttt{setText}, ecc.).

\textbf{Jetpack Compose} sposa il paradigma \textbf{Dichiarativo}. L'interfaccia utente è scritta interamente in Kotlin tramite funzioni. Lo sviluppatore descrive \textit{cosa} la UI deve mostrare in base allo stato attuale, e Compose si occupa di aggiornarla automaticamente quando lo stato cambia.

\subsection{Regole base di Compose}
\begin{itemize}
    \item \textbf{Annotazione \texttt{@Composable}:} Qualsiasi funzione che emette UI deve essere annotata con \texttt{@Composable}.
    \item \textbf{PascalCase:} Per convenzione, le funzioni composable che fungono da widget si scrivono con l'iniziale maiuscola (es. \texttt{Greeting()}, \texttt{HomeScreen()}), comportandosi di fatto come se fossero delle classi grafiche.
    \item \textbf{Nessun XML:} Tutta la UI risiede in file \texttt{.kt}.
\end{itemize}

\section{Componenti Core e Layout}
I layout classici in XML hanno equivalenti diretti in Compose:

\begin{table}[h]
    \centering
    \begin{tabularx}{\textwidth}{|X|X|}
\hline
\rowcolor{gray!20} \textbf{Mondo XML (Imperativo)} & \textbf{Mondo Compose (Dichiarativo)} \\ \hline
\texttt{LinearLayout (vertical)} & \texttt{Column \{ \}} \\ \hline
\texttt{LinearLayout (horizontal)} & \texttt{Row \{ \}} \\ \hline
\texttt{FrameLayout} (Sovrapposizione) & \texttt{Box \{ \}} \\ \hline
\texttt{TextView} & \texttt{Text("Stringa")} \\ \hline
\texttt{RecyclerView} / \texttt{ListView} & \texttt{LazyColumn \{ \}} / \texttt{LazyRow \{ \}} \\ \hline
\texttt{setOnClickListener \{ \}} & \texttt{Modifier.clickable \{ \}} oppure \texttt{Button(onClick=...) \{ \}}\\ \hline
    \end{tabularx}
\end{table}

\subsection{I Modifier}
In Compose non esistono attributi XML come \texttt{android:layout\_width}. Tutto ciò che riguarda lo stile, le dimensioni, il padding e i comportamenti (click, scroll) viene passato tramite il parametro \textbf{\texttt{Modifier}}.

\begin{lstlisting}[caption={Esempio pratico di Modifier}, label={lst:modifier}]
Text(
    text = "Hello PDM!",
    modifier = Modifier
        .fillMaxWidth()        // Equivalente a match_parent
        .padding(16.dp)        // Margine/Padding
        .clickable { /* fai qualcosa */ }
)
\end{lstlisting}

\section{Lo Stato e la Ricomposizione (Cruciale per l'esame)}
In Compose, \textbf{la UI è una funzione dello Stato: UI = f(state)}.

\begin{center}
    \begin{tikzpicture}[
        node distance=2cm,
        box/.style={rectangle, draw, rounded corners, fill=white, text width=2.5cm, align=center, drop shadow},
        arrow/.style={-Stealth, thick}
    ]
        \node[box, fill=blue!10] (state) {\textbf{State}\\(es. variabile count)};
        \node[box, fill=green!10, below=of state] (ui) {\textbf{UI}\\(es. Text(count))};
        \node[box, fill=orange!10, right=of ui] (event) {\textbf{Event}\\(es. Button Click)};

        \draw[arrow, blue] (state) -- node[left] {Aggiorna} (ui);
        \draw[arrow, orange] (ui) -- node[below] {Genera} (event);
        \draw[arrow, red] (event) |- node[right] {Modifica} (state);
    \end{tikzpicture}
\end{center}

Se una semplice variabile Kotlin (es. \texttt{var count = 0}) cambia, la UI \textbf{non} si aggiorna. Per far capire a Compose che deve ri-disegnare la schermata, dobbiamo usare \textbf{\texttt{State}} e \textbf{\texttt{remember}}.

\begin{itemize}
    \item \textbf{\texttt{mutableStateOf(valore)}:} Crea un contenitore osservabile. Quando il valore cambia, Compose viene "notificato".
    \item \textbf{\texttt{remember \{ \}}:} Dice a Compose di "ricordare" questo valore tra una ricomposizione e l'altra, altrimenti ad ogni refresh la variabile verrebbe re-inizializzata al valore di partenza.
    \item \textbf{Ricomposizione (Recomposition):} È il processo con cui Compose ri-esegue le funzioni \texttt{@Composable} i cui dati di stato sono cambiati, aggiornando solo i pezzi di schermo necessari.
\end{itemize}

\begin{lstlisting}[caption={Gestione corretta dello Stato}, label={lst:compose_state}]
@Composable
fun Counter() {
    // 1. Dichiariamo lo stato
    var count by remember { mutableStateOf(0) }

    Button(onClick = {
        // 2. Modifichiamo lo stato (scatena la Ricomposizione)
        count++
    }) {
        // 3. La UI si aggiorna automaticamente
        Text("Click: $count")
    }
}
\end{lstlisting}

\section{Le Liste Dinamiche: LazyColumn}
Mentre nel vecchio sistema serviva creare una \texttt{RecyclerView}, un \texttt{Adapter}, un \texttt{ViewHolder} e un layout XML per la riga, in Compose le liste sono banali.

\textbf{\texttt{LazyColumn}} (e \texttt{LazyRow}) si chiamano "Lazy" perché instanziano (disegnano) solo gli elementi visibili sullo schermo in quel momento, offrendo grandissime performance.

\begin{lstlisting}[caption={Creare una lista in Compose}, label={lst:lazycolumn}]
@Composable
fun UserListScreen(users: List<String>) {
    LazyColumn {
        // Il blocco 'items' sostituisce il vecchio Adapter!
        items(users) { user ->
            // Qui definiamo la singola riga
            Text(text = "Utente: $user", modifier = Modifier.padding(8.dp))
        }
    }
}
\end{lstlisting}

\section{Esecuzione Asincrona e Immagini (Side Effects)}

\subsection{LaunchedEffect}
Se dobbiamo eseguire codice asincrono (come scaricare un JSON o dati da un database) non possiamo scriverlo direttamente nel corpo di un \texttt{@Composable}, altrimenti verrebbe rieseguito (scaricando ripetutamente i dati) ad ogni singola ricomposizione!

Per evitare questo si usa \textbf{\texttt{LaunchedEffect}}. Viene eseguito una sola volta quando il composable entra nello schermo (o quando una "key" specificata cambia). Per operazioni di rete si usa il blocco \texttt{withContext(Dispatchers.IO)}.

\subsection{Caricamento Immagini da URL}
In Jetpack Compose, per caricare un'immagine da Internet si consiglia la libreria di terze parti \textbf{Coil}, che offre il widget \textbf{\texttt{AsyncImage}}. Fa tutto da solo: apre un thread asincrono, scarica, salva in cache e mostra l'immagine, senza bloccare la UI.

\begin{lstlisting}[caption={Esempio di AsyncImage}, label={lst:asyncimage}]
AsyncImage(
    model = "https://picsum.photos/id/10/200/300",
    contentDescription = "Descrizione per accessibilita",
    modifier = Modifier.size(64.dp)
)
\end{lstlisting}

\section{Navigazione in Compose (Navigation Compose)}
Il Navigation Component è stato adattato per funzionare senza XML.
Servono 3 elementi:
\begin{enumerate}
    \item \textbf{\texttt{rememberNavController()}:} Per tenere traccia dello stato della navigazione.
    \item \textbf{\texttt{NavHost}:} Il contenitore che decide quale Composable mostrare. Sostituisce il NavGraph XML specificando una \texttt{startDestination}.
    \item \textbf{\texttt{composable("rotta")}:} Funzione che associa una stringa (la rotta) al Widget da disegnare.
\end{enumerate}

\begin{lstlisting}[caption={Navigazione in Compose}, label={lst:nav_compose}]
@Composable
fun AppNavigation() {
    val navController = rememberNavController()

    NavHost(navController = navController, startDestination = "home") {
        composable("home") {
            HomeScreen(onNavigate = { navController.navigate("detail") })
        }
        composable("detail") {
            DetailScreen()
        }
    }
}
\end{lstlisting}